\documentclass{zhvt-classic}

% 使用本地源流明体（GenRyuMin2 TW）按文件路径加载，不依赖系统字体名
% 版框（内框／界栏／书口侧栏／鱼尾）改用与版心印章同一印泥棕色：
%   印章图片 yshs.png 实色 = RGB(135,68,52) = HTML 874434
\usepackage{xcolor}
\usepackage{tikz}
% sealbrown 双注册：xcolor 供 TikZ/pgf，L3 供 cls 版框
\definecolor{sealbrown}{HTML}{874434}
\ExplSyntaxOn
\color_set:nnn { sealbrown } { HTML } { 874434 }
\ExplSyntaxOff

\zhvtset{
  grid color       = sealbrown,
  main font        = GenRyuMin2TW-SB,
  main font path   = C:/Users/Administrator/texmf/fonts/opentype/public/GenRyuMin2TW/,
jiazhu font      = GenRyuMin2TW-M,
  jiazhu font path = C:/Users/Administrator/texmf/fonts/opentype/public/GenRyuMin2TW/,
  % 页码框容纳 4 字：正常页「九十七」+前置空 1 格=4 行恰好；
  % 100+ 页「一百零一」共 5 行超出，末字「一/二/三」被框裁掉（2026-10-04 实测），故放宽到 5
  shukou page number chars = 5,
  % 书口书名框高从默认 6 字减到 4 字，把省出的 2 字高拨给 chapter 框，
  % 使 6 字书名+卷名放进 9 字高的 chapter 框后不再压线（2026-10-05）
  shukou title chars = 4,
}

\title{史记}[鲁周公世家]
\author{汉·司马迁}
\maker{闲叟漫记·重製}


% ===== 索隐述赞样式（2026-11-08 四次定稿 · 愚大人亲手改定参数）=====
% ⚠ 竖排前提：自定义字体家族必须带 RawFeature={vertical:+vert}；
%   simhei.ttf 加 +vhal 会整体加载失败静默回退，只写 +vert；simfang 同理。
\newCJKfontfamily{\shuzanfont}[Path=C:/Windows/Fonts/,RawFeature={vertical:+vert}]{simhei.ttf}
\newCJKfontfamily{\shuzanbody}[Path=C:/Windows/Fonts/,RawFeature={vertical:+vert}]{simfang.ttf}
% 题头：152x30pt 圆角 6pt 棕底白字；baseline (0,2) 使整框落入界行格子不压版框线；
% 白字节点 x 不补偿、y +8.2pt（像素级墨迹实测 dx/dy≈0，L/R 留白约 20pt）。
\newcommand{\shuzanhead}[1]{%
  \begin{tikzpicture}[baseline={(0,2)}]
    \fill[sealbrown, rounded corners=6pt]
      (-76pt,-15pt) rectangle (76pt,15pt);
    \node[text=white, inner sep=0pt, anchor=center] at (0pt,8.2pt)
      {\shuzanfont\fontsize{23}{23}\selectfont #1};
  \end{tikzpicture}%
}
% ★ 修正书口卷名位置：chapter 框起始位置改回 cls 默认值，避免全局影响。
%   只在每个卷首页通过 \juan 命令里的空命令微调该页位置（安全，不会跨页污染）。
\ExplSyntaxOn
\AtBeginDocument{
  \dim_set:Nn \l__zhvt_shukou_chapter_title_start_dim
    { \l__zhvt_shukou_title_height_dim + \l__zhvt_font_size_dim * 2 }
  \dim_set:Nn \l__zhvt_shukou_page_number_start_dim
    { \l__zhvt_shukou_chapter_title_start_dim + \l__zhvt_shukou_chapter_title_height_dim }
}
\ExplSyntaxOff

\begin{document}

% ★★ 版心（书口）排布覆写：把书名放到鱼尾之下（2026-10-04）
%
%   cls 默认版心自上而下：书名（title 静态框，在鱼尾之上）→ 鱼尾 → 篇名（chapter 动态框）→ 页码
%   本模板改为：          鱼尾（位置不动）→ 书名 → 篇名 → 页码分割线 → 页码
%
%   手法：把 cls 放在 title 静态框里的书名，改放进 chapter 动态框（鱼尾之下）；
%         title 静态框留空。书名取 \l__zhvt_title_str，即 \title{...} 的第一参数，
%         所以换书时只改 \title，不必在此重复写书名。
%
%   ⚠ 三条必须遵守的写法（2026-10-04 实测）：
%     ① 汉字必须用 \hbox{...} 包住。flowfram 会把框内容装进一个宽 14.4mm（= 一个 line_sep
%        = 书口栏宽）的纵向 minipage 并按该宽度折行，而整页靠 PDF /Rotate 90 显示、
%        框内文字本身不变向 —— 不包 \hbox 时「围炉夜话」4 字会被逐字折成 4 行，
%        在版面上就成了横排 4 字并横向溢出版框。
%     ② 竖排下**一个 \hbox 就是一整列**。书名与篇名要同处一列，必须放进同一个 \hbox；
%        分成两个 \hbox 时，后一个会被当成新的一列、甩到隔壁栏去
%        （实测「序」跑到书口栏外侧 14.4mm 的正文列上）。中间的 \quad 在框内是横向空白，
%        显示出来正好是书名与篇名之间的纵向间隔。
%     ③ ★ 内容盒要写成 **与框等宽的定宽盒** `\hbox to \l__zhvt_shukou_chapter_title_height_dim{…\hfil}`。
%        实测：chapter 框是把内容盒沿纵向**居中**摆放的，内容宽度一变，书名就跟着漂：
%        自然宽度的 \hbox（书名 38.4mm）落在 x=157.2mm，加篇名后（57.6mm）落在 x=152.4mm，
%        目录页（无篇名）与正文页的书名就差半格。做成等宽盒后，无论有没有篇名，
%        内容固定从框顶起排，实测两页均落在 x=142.8mm，不再漂。
\ExplSyntaxOn
% 把书口 chapter 框（书名+篇名）起始位置下移 1 字高，避免 4 字卷名顶到鱼尾
% ⚠ cls 的 \__zhvt_calculate_dimensions: 在 \AtBeginDocument 里执行，会重算该 dim；
%   必须把覆写挂到 begin document 之后，否则会被二次覆盖（2026-10-05 实测）。
% (旧方案已移到导言区，直接修改 chapter 框起始位置)
% 卷标题命令 #1=书口篇名(鱼尾) #2=正文卷名（含卷序，卷名内空格用 ~）
%
% ★★ 目录页两项修订（2026-10-05，愚大人指出）★★
%   ① 页码起点统一对齐：
%      原状 —— 目录项写成 \par\hspace{4em}卷名\hspace{4em}页码，页码起点 = 卷名实宽 + 4em。
%      卷名实宽随字数变：「卷一~梁惠王上」6 字 +~、「卷七~离娄上」5 字 +~，
%      差 1 个字位 → 卷七~卷十 4 条的页码整体比其余 10 条高 27.21pt（一个字），看着上下不齐。
%      改法 —— 卷名补全角空格（宽 = 1 字 = 27.215pt）补足到 **7 字符位**再写进目录，
%      14 条页码起点遂全部落在同一条竖线上（起点对齐；末字随页码字数自然参差）。
%      ⚠ 只对目录里的副本 \l_tmpc_tl 补，正文卷名仍用 \l_tmpa_tl 原样 —— 正文像素不受影响。
%   ② 卷名整体上移半字位：前导 \hspace 由 4em 改 3.5em。
%      实测原状目录块在版框内偏下（天头留白 139.6pt / 地脚留白 110.7pt），
%      上移 0.5em = 13.6pt 后上下基本均分，视觉居中。
\NewDocumentCommand { \juan } { mm } {
	\clearpage
	\tl_set:Nn \l__zhvt_chapter_title_tl { #1 }
	\tl_set:Nn \l_tmpa_tl { #2 }
	\tl_set_eq:NN \l_tmpc_tl \l_tmpa_tl
	\int_step_inline:nn { 7 - \str_count:n { #2 } }
	  { \tl_put_right:Nn \l_tmpc_tl { 　 } }
	\tl_set:Nx \l_tmpb_tl { \zhnum { page } }
	\tl_put_left:Nn \l_tmpb_tl { \hspace { 4em } }
	\tl_put_left:Nn \l_tmpb_tl { \l_tmpc_tl }
	\tl_put_left:Nn \l_tmpb_tl { \par \hspace { 3.5em } }
	\addtocontents { toc } { \l_tmpb_tl }
	\skip_horizontal:n { \l__zhvt_font_size_dim * 4 }
	\l_tmpa_tl
	\par
	\thispagestyle { empty }
}
\cs_set_protected:Npn \__zhvt_set_shukou_content: {
	% ★ 鱼尾上方留白（title 静态框）嵌入闲叟钤印：
	%   整页 /Rotate 90，故此处把图用 \rotatebox{90} 转正为竖章；
	%   尺寸卡在版心条幅内（书内约 7mm 宽 × 33.8mm 高，留余量不压鱼尾）。
	\setstaticcontents * { title } {
		\vbox_to_ht:nn { \l__zhvt_baseline_skip_dim } {
			\vskip 14.2pt
			\hbox_to_wd:nn { \l__zhvt_shukou_title_height_dim } {
				\hfil
				\rotatebox [ origin = c ] { 90 } { \includegraphics [ width = 7mm , height = 33.8mm ] { yshs.png } }
				\hfil
			}
		}
	}
% ★ 版心（书口）统一改法（2026-11-16 愚大人定）：18pt＋框内居中＋书口文字同版框色
%   ① 原 27.2pt 时 7 字长篇名（如「廉颇蔺相如列传」）末字溢出 chapter 框、压住页码；
%      改 \fontsize{18pt}{18pt}，前后 \hfil 令字串在框内居中（不随篇名字数漂移）。
%   ② \vbox{\vskip 8.2pt} 定量下移，抵消缩字号后栏宽方向 8.2pt 上偏（像素实测 y 中心 543.3 = 原版）。
%      ⚠ \raisebox / \vspace* 在 flowfram 框内无效；换字号须同步换补偿：16pt→10.1、19pt→7.4。
%   ③ chapter 与 page~number 均置 \color{sealbrown}（HTML 874434），与版框／界行同色；正文仍黑。
%   ⚠ 页码字位、钤印、鱼尾、界行、正文一律未动；五部合订本同此口径。
	\setdynamiccontents * { chapter }
	{ \color{sealbrown} \vbox { \vskip 8.2pt \hbox to \l__zhvt_shukou_chapter_title_height_dim
		{ \hfil \fontsize{18pt}{18pt}\selectfont \l__zhvt_title_str \quad \l__zhvt_chapter_title_tl \hfil } } }
	\setdynamiccontents * { page~number } { \color{sealbrown} \quad \zhnum { page } }
}
\__zhvt_set_shukou_content:
\ExplSyntaxOff

\maketitle

% 单篇，无多卷目录，不渲染目录页

\juan{鲁周公世家}{卷三十三~鲁周公世家第三}
\hspace*{2em}周公旦者。周武王弟也。自文王在时。旦为子孝。笃仁。异于群子。及武王即位。旦常辅翼武王。用事居多。武王九年。东伐至盟津。周公辅行。十一年。伐纣。至牧野。周公佐武王。作牧誓。破殷。入商宫。已杀纣。周公把大钺。召公把小钺。以夹武王。衅社。告纣之罪于天。及殷民。释箕子之囚。封纣子武庚禄父。使管叔。蔡叔傅之。以续殷祀。遍封功臣同姓戚者。封周公旦于少昊之虚曲阜。是为鲁公。周公不就封。留佐武王。\\
\hspace*{2em}武王克殷二年。天下未集。武王有疾。不豫。群臣惧。太公。召公乃缪卜。周公曰。﹁未可以戚我先王。﹂周公于是乃自以为质。设三坛。周公北面立。戴璧秉圭。告于太王。王季。文王。史策祝曰。﹁惟尔元孙王发。勤劳阻疾。若尔三王是有负子之责于天。以旦代王发之身。旦巧能。多材多艺。能事鬼神。乃王发不如旦多材多艺。不能事鬼神。乃命于帝庭。敷佑四方。用能定汝子孙于下地。四方之民罔不敬畏。无坠天之降葆命。我先王亦永有所依归。今我其即命于元龟。尔之许我。我以其璧与圭归。以俟尔命。尔不许我。我乃屏璧与圭。﹂周公已令史策告太王。王季。文王。欲代武王发。于是乃即三王而卜。卜人皆曰吉。发书视之。信吉。周公喜。开籥。乃见书遇吉。周公入贺武王曰。﹁王其无害。旦新受命三王。维长终是图。兹道能念予一人。﹂周公藏其策金縢匮中。诫守者勿敢言。明日。武王有瘳。\\
\hspace*{2em}其后武王既崩。成王少。在强葆之中。周公恐天下闻武王崩而畔。周公乃践阼代成王摄行政当国。管叔及其群弟流言于国曰。﹁周公将不利于成王。﹂周公乃告太公望。召公奭曰。﹁我之所以弗辟而摄行政者。恐天下畔周。无以告我先王太王。王季。文王。三王之忧劳天下久矣。于今而后成。武王蚤终。成王少。将以成周。我所以为之若此。﹂于是卒相成王。而使其子伯禽代就封于鲁。周公戒伯禽曰。﹁我文王之子。武王之弟。成王之叔父。我于天下亦不贱矣。然我一沐三捉发。一饭三吐哺。起以待士。犹恐失天下之贤人。子之鲁。慎无以国骄人。﹂\\
\hspace*{2em}管。蔡。武庚等果率淮夷而反。周公乃奉成王命。兴师东伐。作大诰。遂诛管叔。杀武庚。放蔡叔。收殷余民。以封康叔于卫。封微子于宋。以奉殷祀。宁淮夷东土。二年而毕定。诸侯咸服宗周。\\
\hspace*{2em}天降祉福。唐叔得禾。异母同颖。献之成王。成王命唐叔以馈周公于东土。作馈禾。周公既受命禾。嘉天子命。作嘉禾。东土以集。周公归报成王。乃为诗贻王。命之曰鸱鸮。王亦未敢训周公。\\
\hspace*{2em}成王七年二月乙未。王朝步自周。至丰。使太保召公先之雒相土。其三月。周公往营成周雒邑。卜居焉。曰吉。遂国之。\\
\hspace*{2em}成王长。能听政。于是周公乃还政于成王。成王临朝。周公之代成王治。南面倍依以朝诸侯。及七年后。还政成王。北面就臣位。匔匔如畏然。\\
\hspace*{2em}初。成王少时。病。周公乃自揃其蚤沉之河。以祝于神曰。﹁王少未有识。奸神命者乃旦也。﹂亦藏其策于府。成王病有瘳。及成王用事。人或谮周公。周公奔楚。成王发府。见周公祷书。乃泣。反周公。\\
\hspace*{2em}周公归。恐成王壮。治有所淫佚。乃作多士。作毋逸。毋逸称。﹁为人父母。为业至长久。子孙骄奢忘之。以亡其家。为人子可不慎乎。故昔在殷王中宗。严恭敬畏天命。自度治民。震惧不敢荒宁。故中宗飨国七十五年。其在高宗。久劳于外。为与小人。作其即位。乃有亮闇。三年不言。言乃欢。不敢荒宁。密靖殷国。至于小大无怨。故高宗飨国五十五年。其在祖甲。不义惟王。久为小人于外。知小人之依。能保施小民。不侮鳏寡。故祖甲飨国三十三年。﹂多士称曰。﹁自汤至于帝乙。无不率祀明德。帝无不配天者。在今后嗣王纣。诞淫厥佚。不顾天及民之从也。其民皆可诛。﹂﹁文王日中昃不暇食。飨国五十年。﹂作此以诫成王。\\
\hspace*{2em}成王在丰。天下已安。周之官政未次序。于是周公作周官。官别其宜。作立政。以便百姓。百姓说。\\
\hspace*{2em}周公在丰。病。将没。曰。﹁必葬我成周。以明吾不敢离成王。﹂周公既卒。成王亦让。葬周公于毕。从文王。以明予小子不敢臣周公也。\\
\hspace*{2em}周公卒后。秋未获。暴风雷。禾尽偃。大木尽拔。周国大恐。成王与大夫朝服以开金縢书。王乃得周公所自以为功代武王之说。二公及王乃问史百执事。史百执事曰。﹁信有。昔周公命我勿敢言。﹂成王执书以泣。曰。﹁自今后其无缪卜乎。昔周公勤劳王家。惟予幼人弗及知。今天动威以彰周公之德。惟朕小子其迎。我国家礼亦宜之。﹂王出郊。天乃雨。反风。禾尽起。二公命国人。凡大木所偃。尽起而筑之。岁则大孰。于是成王乃命鲁得郊祭文王。鲁有天子礼乐者。以褒周公之德也。\\
\hspace*{2em}周公卒。子伯禽固已前受封。是为鲁公。鲁公伯禽之初受封之鲁。三年而后报政周公。周公曰。﹁何迟也。﹂伯禽曰。﹁变其俗。革其礼。丧三年然后除之。故迟。﹂太公亦封于齐。五月而报政周公。周公曰。﹁何疾也。﹂曰。﹁吾简其君臣礼。从其俗为也。﹂及后闻伯禽报政迟。乃叹曰。﹁呜呼。鲁后世其北面事齐矣。夫政不简不易。民不有近。平易近民。民必归之。﹂\\
\hspace*{2em}伯禽即位之后。有管。蔡等反也。淮夷。徐戎亦并兴反。于是伯禽率师伐之于肸。作肸誓。曰。﹁陈尔甲胄。无敢不善。无敢伤牿。马牛其风。臣妾逋逃。勿敢越逐。敬复之。无敢寇攘。逾墙垣。鲁人三郊三隧。歭尔刍茭。糗粮。桢干。无敢不逮。我甲戌筑而征徐戎。无敢不及。有大刑。﹂作此肸誓。遂平徐戎。定鲁。\\
\hspace*{2em}鲁公伯禽卒。子考公酋立。考公四年卒。立弟熙。是谓炀公。炀公筑茅阙门。六年卒。子幽公宰立。幽公十四年。幽公弟杀幽公而自立。是为魏公。魏公五十年卒。子厉公擢立。厉公三十七年卒。鲁人立其弟具。是为献公。献公三十二年卒。子真公濞立。\\
\hspace*{2em}真公十四年。周厉王无道。出奔彘。共和行政。二十九年。周宣王即位。\\
\hspace*{2em}三十年。真公卒。弟敖立。是为武公。\\
\hspace*{2em}武公九年春。武公与长子括。少子戏。西朝周宣王。宣王爱戏。欲立戏为鲁太子。周之樊仲山父谏宣王曰。﹁废长立少。不顺。不顺。必犯王命。犯王命。必诛之。故出令不可不顺也。令之不行。政之不立。行而不顺。民将弃上。夫下事上。少事长。所以为顺。今天子建诸侯。立其少。是教民逆也。若鲁从之。诸侯效之。王命将有所壅。若弗从而诛之。是自诛王命也。诛之亦失。不诛亦失。王其图之。﹂宣王弗听。卒立戏为鲁太子。夏。武公归而卒。戏立。是为懿公。\\
\hspace*{2em}懿公九年。懿公兄括之子伯御与鲁人攻弑懿公。而立伯御为君。伯御即位十一年。周宣王伐鲁。杀其君伯御。而问鲁公子能道顺诸侯者。以为鲁后。樊穆仲曰。﹁鲁懿公弟称。肃恭明神。敬事耆老。赋事行刑。必问于遗训而咨于固实。不干所问。不犯所咨。﹂宣王曰。﹁然。能训治其民矣。﹂乃立称于夷宫。是为孝公。自是后。诸侯多畔王命。\\
\hspace*{2em}孝公二十五年。诸侯畔周。犬戎杀幽王。秦始列为诸侯。\\
\hspace*{2em}二十七年。孝公卒。子弗湟立。是为惠公。\\
\hspace*{2em}惠公三十年。晋人弑其君昭侯。四十五年。晋人又弑其君孝侯。\\
\hspace*{2em}四十六年。惠公卒。长庶子息摄当国。行君事。是为隐公。初。惠公適夫人无子。公贱妾声子生子息。息长。为娶于宋。宋女至而好。惠公夺而自妻之。生子允。登宋女为夫人。以允为太子。及惠公卒。为允少故。鲁人共令息摄政。不言即位。\\
\hspace*{2em}隐公五年。观渔于棠。八年。与郑易天子之太山之邑祊及许田。君子讥之。\\
\hspace*{2em}十一年冬。公子挥谄谓隐公曰。﹁百姓便君。君其遂立。吾请为君杀子允。君以我为相。﹂隐公曰。﹁有先君命。吾为允少。故摄代。今允长矣。吾方营菟裘之地而老焉。以授子允政。﹂挥惧子允闻而反诛之。乃反谮隐公于子允曰。﹁隐公欲遂立。去子。子其图之。请为子杀隐公。﹂子允许诺。十一月。隐公祭钟巫。齐于社圃。馆于蒍氏。挥使人弑隐公于蒍氏。而立子允为君。是为桓公。\\
\hspace*{2em}桓公元年。郑以璧易天子之许田。二年。以宋之赂鼎入于太庙。君子讥之。三年。使挥迎妇于齐为夫人。六年。夫人生子。与桓公同日。故名曰同。同长。为太子。\\
\hspace*{2em}十六年。会于曹。伐郑。入厉公。\\
\hspace*{2em}十八年春。公将有行。遂与夫人如齐。申繻谏止。公不听。遂如齐。齐襄公通桓公夫人。公怒夫人。夫人以告齐侯。夏四月丙子。齐襄公飨公。公醉。使公子彭生抱鲁桓公。因命彭生折其胁。公死于车。鲁人告于齐曰。﹁寡君畏君之威。不敢宁居。来修好礼。礼成而不反。无所归咎。请得彭生以除丑于诸侯。﹂齐人杀彭生以说鲁。立太子同。是为庄公。庄公母夫人因留齐。不敢归鲁。\\
\hspace*{2em}庄公五年冬。伐卫。内卫惠公。\\
\hspace*{2em}八年。齐公子纠来奔。九年。鲁欲内子纠于齐。后桓公。桓公发兵击鲁。鲁急。杀子纠。召忽死。齐告鲁生致管仲。鲁人施伯曰。﹁齐欲得管仲。非杀之也。将用之。用之则为鲁患。不如杀。以其尸与之。﹂庄公不听。遂囚管仲与齐。齐人相管仲。\\
\hspace*{2em}十三年。鲁庄公与曹沬会齐桓公于柯。曹沬劫齐桓公。求鲁侵地。已盟而释桓公。桓公欲背约。管仲谏。卒归鲁侵地。十五年。齐桓公始霸。二十三年。庄公如齐观社。\\
\hspace*{2em}三十二年。初。庄公筑台临党氏。见孟女。说而爱之。许立为夫人。割臂以盟。孟女生子斑。斑长。说梁氏女。往观。圉人荦自墙外与梁氏女戏。斑怒。鞭荦。庄公闻之。曰。﹁荦有力焉。遂杀之。是未可鞭而置也。﹂斑未得杀。会庄公有疾。庄公有三弟。长曰庆父。次曰叔牙。次曰季友。庄公取齐女为夫人曰哀姜。哀姜无子。哀姜娣曰叔姜。生子开。庄公无适嗣。爱孟女。欲立其子斑。庄公病。而问嗣于弟叔牙。叔牙曰。﹁一继一及。鲁之常也。庆父在。可为嗣。君何忧。﹂庄公患叔牙欲立庆父。退而问季友。季友曰。﹁请以死立斑也。﹂庄公曰。﹁曩者叔牙欲立庆父。奈何。﹂季友以庄公命命牙待于针巫氏。使针季劫饮叔牙以鸩。曰。﹁饮此则有后奉祀。不然。死且无后。﹂牙遂饮鸩而死。鲁立其子为叔孙氏。八月癸亥。庄公卒。季友竟立子斑为君。如庄公命。侍丧。舍于党氏。\\
\hspace*{2em}先时庆父与哀姜私通。欲立哀姜娣子开。及庄公卒而季友立斑。十月己未。庆父使圉人荦杀鲁公子斑于党氏。季友奔陈。庆父竟立庄公子开。是为愍公。\\
\hspace*{2em}愍公二年。庆父与哀姜通益甚。哀姜与庆父谋杀愍公而立庆父。庆父使卜齮袭杀愍公于武闱。季友闻之。自陈与愍公弟申如邾。请鲁求内之。鲁人欲诛庆父。\\
\hspace*{2em}庆父恐。奔莒。于是季友奉子申入。立之。是为釐公。釐公亦庄公少子。哀姜恐。奔邾。季友以赂如莒求庆父。庆父归。使人杀庆父。庆父请奔。弗听。乃使大夫傒斯行哭而往。庆父闻傒斯音。乃自杀。齐桓公闻哀姜与庆父乱以危鲁。及召之邾而杀之。以其尸归。戮之鲁。鲁釐公请而葬之。\\
\hspace*{2em}季友母陈女。故亡在陈。陈故佐送季友及子申。季友之将生也。父鲁桓公使人卜之。曰。﹁男也。其名曰﹃友﹄。间于两社。为公室辅。季友亡。则鲁不昌。﹂及生。有文在掌曰﹁友﹂。遂以名之。号为成季。其后为季氏。庆父后为孟氏也。\\
\hspace*{2em}釐公元年。以汶阳鄪封季友。季友为相。\\
\hspace*{2em}九年。晋里克杀其君傒齐。卓子。齐桓公率釐公讨晋乱。至高梁而还。立晋惠公。十七年。齐桓公卒。二十四年。晋文公即位。\\
\hspace*{2em}三十三年。釐公卒。子兴立。是为文公。\\
\hspace*{2em}文公元年。楚太子商臣弑其父成王。代立。三年。文公朝晋襄公。\\
\hspace*{2em}十一年十月甲午。鲁败翟于咸。获长翟乔如。富父终甥舂其喉以戈。杀之。埋其首于子驹之门。以命宣伯。\\
\hspace*{2em}初。宋武公之世。鄋瞒伐宋。司徒皇父帅师御之。以败翟于长丘。获长翟缘斯。晋之灭路。获乔如弟棼如。齐惠公二年。鄋瞒伐齐。齐王子城父获其弟荣如。埋其首于北门。卫人获其季弟简如。鄋瞒由是遂亡。\\
\hspace*{2em}十五年。季文子使于晋。\\
\hspace*{2em}十八年二月。文公卒。文公有二妃。长妃齐女为哀姜。生子恶及视。次妃敬嬴。嬖爱。生子俀。俀私事襄仲。襄仲欲立之。叔仲曰不可。襄仲请齐惠公。惠公新立。欲亲鲁。许之。冬十月。襄仲杀子恶及视而立俀。是为宣公。哀姜归齐。哭而过市。曰。﹁天乎。襄仲为不道。杀適立庶。﹂市人皆哭。鲁人谓之﹁哀姜﹂。鲁由此公室卑。三桓强。\\
\hspace*{2em}宣公俀十二年。楚庄王强。围郑。郑伯降。复国之。\\
\hspace*{2em}十八年。宣公卒。子成公黑肱立。是为成公。季文子曰。﹁使我杀適立庶失大援者。襄仲。﹂襄仲立宣公。公孙归父有宠。宣公欲去三桓。与晋谋伐三桓。会宣公卒。季文子怨之。归父奔齐。\\
\hspace*{2em}成公二年春。齐伐取我隆。夏。公与晋郄克败齐顷公于鞍。齐复归我侵地。四年。成公如晋。晋景公不敬鲁。鲁欲背晋合于楚。或谏。乃不。十年。成公如晋。晋景公卒。因留成公送葬。鲁讳之。十五年。始与吴王寿梦会钟离。\\
\hspace*{2em}十六年。宣伯告晋。欲诛季文子。文子有义。晋人弗许。\\
\hspace*{2em}十八年。成公卒。子午立。是为襄公。是时襄公三岁也。\\
\hspace*{2em}襄公元年。晋立悼公。往年冬。晋栾书弑其君厉公。四年。襄公朝晋。\\
\hspace*{2em}五年。季文子卒。家无衣帛之妾。厩无食粟之马。府无金玉。以相三君。君子曰。﹁季文子廉忠矣。﹂\\
\hspace*{2em}九年。与晋伐郑。晋悼公冠襄公于卫。季武子从。相行礼。\\
\hspace*{2em}十一年。三桓氏分为三军。\\
\hspace*{2em}十二年。朝晋。十六年。晋平公即位。二十一年。朝晋平公。\\
\hspace*{2em}二十二年。孔丘生。\\
\hspace*{2em}二十五年。齐崔杼弑其君庄公。立其弟景公。\\
\hspace*{2em}二十九年。吴延陵季子使鲁。问周乐。尽知其意。鲁人敬焉。\\
\hspace*{2em}三十一年六月。襄公卒。其九月。太子卒。鲁人立齐归之子裯为君。是为昭公。\\
\hspace*{2em}昭公年十九。犹有童心。穆叔不欲立。曰。﹁太子死。有母弟可立。不即立长。年钧择贤。义钧则卜之。今裯非適嗣。且又居丧意不在戚而有喜色。若果立。必为季氏忧。﹂季武子弗听。卒立之。比及葬。三易衰。君子曰。﹁是不终也。﹂\\
\hspace*{2em}昭公三年。朝晋至河。晋平公谢还之。鲁耻焉。四年。楚灵王会诸侯于申。昭公称病不往。七年。季武子卒。八年。楚灵王就章华台。召昭公。昭公往贺。赐昭公宝器。已而悔。复诈取之。十二年。朝晋至河。晋平公谢还之。十三年。楚公子弃疾弑其君灵王。代立。十五年。朝晋。晋留之葬晋昭公。鲁耻之。二十年。齐景公与晏子狩竟。因入鲁问礼。二十一年。朝晋至河。晋谢还之。\\
\hspace*{2em}二十五年春。鸜鹆来巢。师己曰。﹁文成之世童谣曰﹃鸜鹆来巢。公在干侯。鸜鹆入处。公在外野﹄。﹂\\
\hspace*{2em}季氏与郈氏斗鸡。季氏芥鸡羽。郈氏金距。季平子怒而侵郈氏。郈昭伯亦怒平子。臧昭伯之弟会伪谗臧氏。匿季氏。臧昭伯囚季氏人。季平子怒。囚臧氏老。臧。郈氏以难告昭公。昭公九月戊戌伐季氏。遂入。平子登台请曰。﹁君以谗不察臣罪。诛之。请迁沂上。﹂弗许。请囚于鄪。弗许。请以五乘亡。弗许。子家驹曰。﹁君其许之。政自季氏久矣。为徒者众。众将合谋。﹂弗听。郈氏曰。﹁必杀之。﹂叔孙氏之臣戾谓其众曰。﹁无季氏与有。孰利。﹂皆曰。﹁无季氏是无叔孙氏。﹂戾曰。﹁然。救季氏。﹂遂败公师。孟懿子闻叔孙氏胜。亦杀郈昭伯。郈昭伯为公使。故孟氏得之。三家共伐公。公遂奔。己亥。公至于齐。齐景公曰。﹁请致千社待君。﹂子家曰。﹁弃周公之业而臣于齐。可乎。﹂乃止。子家曰。﹁齐景公无信。不如早之晋。﹂弗从。叔孙见公还。见平子。平子顿首。初欲迎昭公。孟孙。季孙后悔。乃止。\\
\hspace*{2em}二十六年春。齐伐鲁。取郓而居昭公焉。夏。齐景公将内公。令无受鲁赂。申丰。汝贾许齐臣高龁。子将粟五千庾。子将言于齐侯曰。﹁群臣不能事鲁君。有异焉。宋元公为鲁如晋。求内之。道卒。叔孙昭子求内其君。无病而死。不知天弃鲁乎。抑鲁君有罪于鬼神也。愿君且待。﹂齐景公从之。\\
\hspace*{2em}二十八年。昭公如晋。求入。季平子私于晋六卿。六卿受季氏赂。谏晋君。晋君乃止。居昭公干侯。二十九年。昭公如郓。齐景公使人赐昭公书。自谓﹁主君﹂。昭公耻之。怒而去干侯。三十一年。晋欲内昭公。召季平子。平子布衣跣行。因六卿谢罪。六卿为言曰。﹁晋欲内昭公。众不从。﹂晋人止。三十二年。昭公卒于干侯。鲁人共立昭公弟宋为君。是为定公。\\
\hspace*{2em}定公立。赵简子问史墨曰。﹁季氏亡乎。﹂史墨对曰。﹁不亡。季友有大功于鲁。受鄪为上卿。至于文子。武子。世增其业。鲁文公卒。东门遂杀適立庶。鲁君于是失国政。政在季氏。于今四君矣。民不知君。何以得国。是以为君慎器与名。不可以假人。﹂\\
\hspace*{2em}定公五年。季平子卒。阳虎私怒。囚季桓子。与盟。乃舍之。七年。齐伐我。取郓。以为鲁阳虎邑以从政。八年。阳虎欲尽杀三桓適。而更立其所善庶子以代之。载季桓子将杀之。桓子诈而得脱。三桓共攻阳虎。阳虎居阳关。九年。鲁伐阳虎。阳虎奔齐。已而奔晋赵氏。\\
\hspace*{2em}十年。定公与齐景公会于夹谷。孔子行相事。齐欲袭鲁君。孔子以礼历阶。诛齐淫乐。齐侯惧。乃止。归鲁侵地而谢过。十二年。使仲由毁三桓城。收其甲兵。孟氏不肯堕城。伐之。不克而止。季桓子受齐女乐。孔子去。\\
\hspace*{2em}十五年。定公卒。子将立。是为哀公。\\
\hspace*{2em}哀公五年。齐景公卒。六年。齐田乞弑其君孺子。\\
\hspace*{2em}七年。吴王夫差强。伐齐。至缯。征百牢于鲁。季康子使子贡说吴王及太宰嚭。以礼诎之。吴王曰。﹁我文身。不足责礼。﹂乃止。\\
\hspace*{2em}八年。吴为邹伐鲁。至城下。盟而去。齐伐我。取三邑。十年。伐齐南边。十一年。齐伐鲁。季氏用冉有有功。思孔子。孔子自卫归鲁。\\
\hspace*{2em}十四年。齐田常弑其君简公于俆州。孔子请伐之。哀公不听。十五年。使子服景伯。子贡为介。適齐。齐归我侵地。田常初相。欲亲诸侯。\\
\hspace*{2em}十六年。孔子卒。\\
\hspace*{2em}二十二年。越王句践灭吴王夫差。\\
\hspace*{2em}二十七年春。季康子卒。夏。哀公患三桓。将欲因诸侯以劫之。三桓亦患公作难。故君臣多间。公游于陵坂。遇孟武伯于街。曰。﹁请问余及死乎。﹂对曰。﹁不知也。﹂公欲以越伐三桓。八月。哀公如陉氏。三桓攻公。公奔于卫。去如邹。遂如越。国人迎哀公复归。卒于有山氏。子宁立。是为悼公。\\
\hspace*{2em}悼公之时。三桓胜。鲁如小侯。卑于三桓之冢。\\
\hspace*{2em}十三年。三晋灭智伯。分其地有之。\\
\hspace*{2em}三十七年。悼公卒。子嘉立。是为元公。元公二十一年卒。子显立。是为穆公。穆公三十三年卒。子奋立。是为共公。共公二十二年卒。子屯立。是为康公。康公九年卒。子匽立。是为景公。景公二十九年卒。子叔立。是为平公。是时六国皆称王。\\
\hspace*{2em}平公十二年。秦惠王卒。二十年。平公卒。子贾立。是为文公。文公元年。楚怀王死于秦。二十三年。文公卒。子雠立。是为顷公。\\
\hspace*{2em}顷公二年。秦拔楚之郢。楚顷王东徙于陈。十九年。楚伐我。取徐州。二十四年。楚考烈王伐灭鲁。顷公亡。迁于下邑。为家人。鲁绝祀。顷公卒于柯。\\
\hspace*{2em}鲁起周公至顷公。凡三十四世。\\
\hspace*{2em}太史公曰。余闻孔子称曰﹁甚矣鲁道之衰也。洙泗之间龂龂如也﹂。观庆父及叔牙闵公之际。何其乱也。隐桓之事。襄仲杀適立庶。三家北面为臣。亲攻昭公。昭公以奔。至其揖让之礼则从矣。而行事何其戾也。\\

% ===== 索隐述赞（司马贞《史记索隐》）· 统一样式 =====
% 题头「索隐述赞曰」五字并入反白框；正文每四句一列，定宽盒 \hfil 居列中；
% \fontsize{24}{41} 列距对齐全书界行（121+41k）；\clearpage 使述赞独占一页。
\clearpage
\shuzanhead{索隐述赞曰}\linebreak
{\shuzanbody\fontsize{24}{41}\selectfont
\hbox to \hsize{\hfil 武王既没。成王幼孤。周公摄政。负扆据图。\hfil}\linebreak
\hbox to \hsize{\hfil 及还臣列。北面躬如。元子封鲁。少昊之墟。\hfil}\linebreak
\hbox to \hsize{\hfil 夹辅王室。系职不渝。降及孝公。穆仲致誉。\hfil}\linebreak
\hbox to \hsize{\hfil 隠能让国。春秋之初。丘明执简。褒贬备书。\hfil}%
\par}
\end{document}
